% 10. Conclusion
\section{Conclusion}

This project addressed the long-standing tension between physical
realism and real-time performance in dense crowd simulation.  Position-
Based Dynamics offers stable, anticipatory collision handling, but its
computational cost is dominated by the local interaction density of the
crowd---precisely the regime in which large-scale simulations are most
interesting.  Existing CLOD strategies for crowd simulation reduce
work by allocating fidelity according to camera distance, an inheritance
from rendering LOD that is poorly suited to the physics layer because
the cost of resolving a contact does not depend on how far away the
contact is.

We built a Unity-based crowd simulator that exposes seven solver
backends sharing a common scene, agent representation, and NavMesh
steering pipeline.  The central contribution is \textit{PBDcmp}, a
GPU compute-shader PBD solver in which physics fidelity is allocated
by a \emph{density-driven} continuous level-of-detail system.  The
density field is read directly from the per-bucket counter of the
spatial hash that the solver already maintains for neighbour search,
so density-aware LOD assignment requires no extra buffer, no extra
kernel pass, and no additional memory traffic.  LOD assignment runs
fully in parallel on the GPU, every agent resolves its constraints on
every frame (eliminating frame-offset stuttering), and asynchronous
GPU readback with version-guarded callbacks keeps the main thread
free of physics stalls.

In the adversarial \textit{Bottleneck} test scene at populations of
10--500 agents, the compute solver produced the highest frame rate
and the lowest solver-only cost at every tested population.  At 500
agents it ran at 1001~FPS---$2.07\times$ faster than the
single-threaded Social Force baseline and $7.0\times$ faster than the
single-threaded PBD baseline.  Comparing the parallel-no-LOD baseline
against the full density-driven CLOD pipeline showed a 38\% reduction
in solver cost and a 43\% improvement in end-to-end frame rate from
LOD alone.  Quality cost was minimal: the spatial fidelity reduction
is invisible at typical viewing distances and self-correcting when
clusters densify, in contrast to the temporal stuttering produced by
camera-distance frame-skipping CLOD.

The main takeaway is that density---rather than camera distance---is
the right signal for allocating physics fidelity in a crowd, and that
the data structure required to make this decision is already paid for
by the spatial hash that any reasonable PBD implementation must
maintain.  Density-driven CLOD is therefore not only effective but
essentially free, and combining it with a GPU compute backend yields a
crowd solver that is both faster and visually smoother than every CPU
alternative we tested.

\subsection{Future Work}
\label{sec:future}

Several directions naturally extend this work:

\begin{itemize}[topsep=2pt,itemsep=2pt]
    \item \textbf{Adaptive thresholds.}  The current density
    thresholds $D_{\mathrm{high}}$ and $D_{\mathrm{med}}$ are fixed
    constants tuned by hand.  Adapting them per-frame to match a
    target solver-cost budget---similar to the adaptive radius
    heuristic in \textit{PBDCLODJobs}---would let the solver shed load
    automatically when the GPU is under pressure from another stage of
    the application.

    \item \textbf{Larger-scale validation.}  All measurements were
    collected at populations of 10--500 agents.  The compute solver's
    headroom (still above 1000~FPS at 500 agents) suggests it should
    scale to several thousand agents on the same hardware, but
    confirming this requires test scenes large enough to host the
    larger populations and a profiler harness that automates the sweep.

    \item \textbf{Wider scene library.}  The \textit{Bottleneck} scene
    deliberately stresses contact density, but it is a single scene.
    Open plazas, branching corridors, and multi-bottleneck layouts
    would each exercise the LOD system differently and would either
    confirm or qualify the present results.

    \item \textbf{Integration with perceptual CLOD.}  Density-driven
    physics LOD is orthogonal to perceptual animation
    LOD~\cite{stancu}.  A hybrid system that uses density to drive
    the solver and gaze direction or eccentricity to drive the
    skinning pipeline would compound the savings of both techniques
    and would be a natural extension for VR or head-mounted display
    deployments where gaze data is available.

    \item \textbf{Group behaviour.}  The current LOD treats every
    agent independently.  Promoting an agent to the \textsc{High} tier
    when one of its socially-linked group members is already in
    \textsc{High} would let the system preserve coherent group
    behaviour even when only part of the group is in a dense pocket.

    \item \textbf{Ablation study.}  The compute solver gains its
    speedup from three compounding sources (GPU dispatch, LOD-induced
    work reduction, and elimination of the serial LOD-assignment
    pass).  An ablation that disables each in turn would quantify
    their individual contributions and guide further optimisation.
\end{itemize}
